\paragraph{Extraction cost.} We timed extraction separately on the otherwise idle A10G host (AMD EPYC 7R32), because the timings recorded during the benchmark include contention from dozens of concurrent workers. Each image method fingerprinted the same 256 registered originals at 1024 pixels on the long side, CPU methods on a single thread and learned methods on the GPU in batches of 64 including the host-to-device copy. The hashes cost between \val{fp:lat/image/BlockMean} (BlockMean) and \val{fp:lat/image/wHash-64}\,ms (wHash) per image on one CPU thread, with PDQ at \val{fp:lat/image/PDQ}\,ms and the ISCC Image-Code at \val{fp:lat/image/ISCC-Image-64}\,ms; the learned descriptors cost \val{fp:lat/image/OpenCLIP-B32}--\val{fp:lat/image/ISC21-1st}\,ms on the GPU (DINOv2-S \val{fp:lat/image/DINOv2-S}, SSCD-R50 \val{fp:lat/image/SSCD-mixup}, ISC21 \val{fp:lat/image/ISC21-1st}). None of these costs matters for a signing service: on the same on-demand instance, fingerprinting a million images costs about US\$\val{fp:lat/image/PDQ/usd1m} with PDQ on its eight vCPUs, US\$\val{fp:lat/image/DINOv2-S/usd1m} with DINOv2-S and US\$\val{fp:lat/image/ISC21-1st/usd1m} with ISC21 on its GPU, before decoding. Storage differs more: PDQ, the ISCC codes and the keyed projection need 8--32 bytes per asset, a float DINOv2-S descriptor 1.5\,kB. Audio fingerprints ran at \val{fp:lat/audio/audfprint}--\val{fp:lat/audio/ISCC-Audio-256}\,ms per second of audio (CLAP \val{fp:lat/audio/CLAP}\,ms on the GPU), and video fingerprints at \val{fp:lat/video/TMK+PDQF}\,ms (TMK+PDQF) to \val{fp:lat/video/vPDQ}\,ms (vPDQ) per second of video on one CPU thread and \val{fp:lat/video/DINOv2-S-seq}\,ms for frame-level DINOv2-S, dominated in every case by decoding and frame sampling. We did not time registry search at scale; exact search over our \val{fp:count/image/nref}-item image index was not the bottleneck at any point.
